\section{The Church's Last Words}

When the sacraments have been given and the agony begins, the Church goes on praying. The Roman Ritual provides an order for the commendation of a departing soul, and the Baltimore \emph{Manual} gives its rubric in English: ``Then while the soul is in the agony of its departure, the Priest recites the following prayers.''\endnote{\emph{Rituale Romanum}, Ordo commendationis animae, rubric before the prayer \latin{Proficiscere}, in the English of \emph{A Manual of Prayers} (1889), 514.} Some of the prayers are addressed to God for the dying man, and some to the dying man himself. The first of them begins \latin{Proficiscere, anima christiana}, ``Go forth, O Christian soul,'' and it is the Church's farewell to her child.

A prayer for pardon follows, \latin{Deus misericors}, and after it the commendation proper, which opens with the words on the debt of humanity already heard and then surrounds the departing soul with the whole court of heaven: ``When, therefore, thy soul goeth forth from thy body, may the glorious company of Angels meet thee; may the council of the Apostles who shall judge the world greet thee; may the triumphant army of white-robed Martyrs come out to welcome thee; may the band of shining Confessors, crowned with lilies, encircle thee; may the choir of joyous Virgins receive thee; and unto the bosom of blessed rest may the embrace of the Patriarchs clasp thee; gentle and joyful may the Face of Jesus Christ appear to thee, and may He award thee a place among those who stand before Him for ever.''\endnote{\emph{Rituale Romanum}, Ordo commendationis animae, the prayers \latin{Deus misericors} and \latin{Commendo te}, in \emph{A Manual of Prayers} (1889), 515--517.}

The prayer also names the danger it asks God to avert, ``the terror of darkness, the gnashing of teeth in the flames, the agonies of torment,'' and prays that the soul may never know them. It bids Satan ``give way'' and ``flee away into the vast chaos of eternal night.'' It ends by asking for the vision of God: ``Mayest thou behold thy Redeemer face to face; and, standing ever before Him, gaze with blessed eyes on the Truth made manifest.''\endnote{\latin{Commendo te}, \emph{A Manual of Prayers} (1889), 517--518.}

As the soul is about to depart, the Baltimore \emph{Manual} directs those present to pray more earnestly than ever, and if the dying can no longer speak, to say again and again in his ear the name of Jesus and such words as these: ``Into Thy hands, O Lord, I commend my spirit.---O Lord Jesus Christ, receive my spirit.---Holy Mary, pray for me.---O Mary, mother of grace, mother of mercy, do thou protect me from the enemy, and receive me at the hour of death.''\endnote{\emph{A Manual of Prayers} (1889), The Last Agony, 531; Luke 23:46.} The first is the prayer with which Christ himself died.

\subsection{Come to his assistance}

The Church marks the moment of death with a responsory, said as soon as the soul has gone:

\begin{quote}
Come to his assistance, ye Saints of God, come forth to meet him, ye Angels of the Lord: Receiving his soul: Offering it in the sight of the Most High.

\emph{V.} May Christ receive thee, who hath called thee, and may the Angels bear thee into Abraham's bosom.

\emph{R.} Receiving, etc.

\emph{V.} Eternal rest grant unto him, O Lord, and let perpetual light shine upon him.

\emph{R.} Offering, etc.\endnote{The responsory \latin{Subvenite, Sancti Dei}, \emph{Rituale Romanum} (1872), Ordo commendationis animae; English from \emph{A Manual of Prayers} (1889), The Last Agony, 531--532.}
\end{quote}

\noindent The earlier prayers spoke to God and to the dying man. In the \latin{Subvenite} the Church turns to the saints and angels, and asks them to come out to meet a soul the living can no longer follow.

\subsection{Life is changed, not taken away}

Then the Church carries her dead into her churches and offers for them the Sacrifice of the Mass. In the Roman Missal of 1920, every Mass for the dead, the Masses of All Souls' Day among them, has the Preface for the Dead:

\begin{quote}
It is truly meet and just, right and profitable for us, at all times, and in all places, to give thanks to Thee, O holy Lord, Father almighty, eternal God: through Christ, our Lord, in Whom the hope of a happy resurrection has shone on us, so that those whom the certain fate of dying renders sad, may be consoled by the promise of future immortality. For with regard to Thy faithful, O Lord, life is changed, not taken away; and the house of their earthly dwelling being destroyed, an eternal dwelling in heaven is obtained.\endnote{The Preface for the Dead, in Lasance, \emph{The New Roman Missal} (1945), Masses for the Dead. Latin in the \emph{Missale Romanum} (typical edition of 1920): \latin{Tuis enim fidelibus, Domine, vita mutatur, non tollitur: et, dissoluta terrestris hujus incolatus domo, \ae terna in coelis habitatio comparatur}.}
\end{quote}

\noindent The Preface admits the sadness before it gives the promise, naming ``the certain fate of dying'' and granting that it ``renders sad.'' That sadness is the violence Augustine would not call gentle, and the promise is the passage to eternity that Cyprian preached in the plague. The image of the destroyed earthly house is the Apostle's.\endnote{2 Corinthians 5:1, Douay--Rheims (Challoner).} The promise itself, \latin{Vita mutatur, non tollitur}, is made to the faithful, whose life is changed by death and not taken away. That life is the one they received in baptism.

The Office of the Dead gives both the sadness and the hope in the words of Job. Its first lesson begs, ``spare me, for my days are nothing.'' Its eighth lesson reads: ``For I know that my Redeemer liveth, and in the last day I shall rise out of the earth. And I shall be clothed again with my skin, and in my flesh I shall see my God.''\endnote{Job 19:25--26, Douay--Rheims (Challoner); the Office of the Dead, Matins, Lectio VIII, \emph{Rituale Romanum} (1872): \latin{Scio enim, quod Redemptor meus vivit, et in novissimo die de terra surrecturus sum}.} Because the body laid in the grave will see God, the Church buries her dead with honour, and she does not cease to pray.
